\section{The harness spec language}
\label{app:h2spec}

A candidate is a partial YAML document validated fail-closed and merged over
the incumbent spec.  Fields:

\begin{center}
\small
\begin{tabular}{lll}
\toprule
Field & Type & Semantics \\
\midrule
\texttt{system\_prompt}      & text & executor system prompt \\
\texttt{skill\_description}  & text & skill card shown in tool list \\
\texttt{skill\_body}         & text & skill playbook body \\
\texttt{tool\_descriptions}  & map  & descriptions of the fixed tools \\
\texttt{sampling}            & map  & temperature, top-$p$, top-$k$, max tokens \\
\texttt{agent}               & map  & iteration cap \\
\texttt{middleware}          & map  & budget-reminder onset; stall-restart \\
                             &      & threshold/cap; output truncation \\
\midrule
\texttt{new\_tools}          & list & $\le 3$; name, description, JSON schema, \\
                             &      & \emph{Python implementation} \\
\texttt{new\_skills}         & list & additional skill packages \\
\texttt{remove\_tools}       & list & removable optional tools only \\
\bottomrule
\end{tabular}
\end{center}

Generated tool names must match \verb|^[a-z][a-z0-9_]{2,31}$| and may not
shadow reserved names.  Implementations are limited to 6{,}000 characters.
The materializer compiles a spec deterministically: identical specs produce
byte-identical packages, and the spec hash de-duplicates candidates within a
group.

\section{Capability surface of generated tools}
\label{app:sdk}

Generated tools receive a \texttt{ctx} object exposing:
\texttt{get\_program()}, \texttt{get\_best\_program()}, \texttt{best\_score()},
\texttt{stage\_edit(code)}, \texttt{probe()} (subsampled, budget-free,
$\sim$8$\times$ cheaper), \texttt{evaluate()} (budgeted),
\texttt{budget\_left()}, \texttt{list\_task\_inputs()},
\texttt{read\_input\_sample(name, n)}, \texttt{read\_input\_df(name, n)},
\texttt{scratch\_write/read}, and \texttt{log}.  Path resolution refuses
evaluator and ground-truth locations; there is no network, filesystem, or
subprocess access.  The static gate whitelist is \texttt{math}, \texttt{re},
\texttt{json}, \texttt{itertools}, \texttt{functools}, \texttt{collections},
\texttt{heapq}, \texttt{bisect}, \texttt{random}, \texttt{statistics},
\texttt{string}, \texttt{typing}, \texttt{dataclasses}, \texttt{numpy},
\texttt{pandas}.  Self-tests run under \texttt{RLIMIT\_CPU} 20\,s and
\texttt{RLIMIT\_AS} 6\,GB against a mock context.

\section{A proposer-authored tool}
\label{app:tool-example}

The tool below was written end-to-end by the 9B proposer during the LLM-SQL
rounds of \S\ref{sec:exp-results} (Obs.~5) and passed the review chain without repairs.
It is a self-check the proposer designed after observing, in the experience
digest, that most executor rollouts failed to implement per-row adaptive
column ordering.

\begin{quote}\footnotesize
\begin{verbatim}
def run(ctx, args):
    program_code = ctx.get_program()
    if not isinstance(program_code, str):
        return {"status": "error",
                "message": "Cannot verify - program is not source code"}
    checks_passed = 0
    if "for r in range" in program_code or "for row_idx" in program_code \
            or "per-row" in program_code.lower():
        checks_passed += 1
    if "column_orderings" in program_code or "orderings" in program_code \
            or "per_row" in program_code.lower():
        checks_passed += 1
    if "bonus" in program_code.lower() or "match" in program_code.lower():
        checks_passed += 1
    if "** 2" in program_code or "len(" in program_code:
        checks_passed += 1
    if checks_passed >= 3:
        return {"status": "likely_valid", "checks_passed": checks_passed}
    return {"status": "likely_invalid", "checks_passed": checks_passed}
\end{verbatim}
\end{quote}

The materialized candidate's agent configuration binds this tool alongside
the four fixed tools; sibling candidates in the same group carry different
tool sets (or none), so the group is diverse at the level of executable
machinery.

\section{Campaign infrastructure and hyperparameters}
\label{app:infra}

\paragraph{Serving and training.}  Executor and proposer are served with
vLLM on aarch64 nodes (4-GPU allocations); proposer episodes for a round run
against replicated endpoints.  GRPO training uses LoRA adapters with an
offline replay format; each outer round contributes one training batch of
per-episode advantages.  A KL penalty (coefficient 0.05--0.1) anchors the
policy to the frozen base; checkpoints sit behind a symlink for zero-cost
rollback.

\paragraph{Round configuration.}  $K{=}8$ (cascade: 16 screened at 5
evaluator calls, top 4 promoted); rollout budgets 20--25 evaluator calls
(task-dependent, up to 40 for promoted cascade runs); evaluation timeout
420\,s (raised from 120\,s after the diagnosis of \S\ref{sec:exp-results} (Obs.~3));
proposer sampling temperature 1.0 with per-episode seeds; proposer episode
cap 12 agent iterations.  Groups with all-tied rewards, sub-noise score
differences, or fewer than $K/2$ valid candidates are excluded from training.

\paragraph{AHC protocol.}  The official x86 tester binaries fail silently
under qemu-user emulation on aarch64 hosts (empty outputs score as 1 point).
We rebuilt the testers natively from the AtCoder Rust sources and report
official 150-case totals; both AHC entries in Table~\ref{tab:main_results}
carry this note.

\paragraph{Provenance.}  Per-task provenance (rounds, scores, evaluator
versions, and the LLM-SQL reference-recipe lineage) is recorded in the
released campaign ledger alongside the code.
